\subsubsection{Réactivité : Technologies pré-lithium}

\setcounter{equation}{0}
\setcounter{Q}{0}

Dans les technologies de stockage électrochimiques pré-lithium, trouver les demi-équations rédox associées à chacun des pôles (en décharge) et proposer une équation de réaction fictive modélisant la réaction. Les équations de réactions seront équilibrées en milieu aqueux.

\Q Accumulateur au plomb : $\ominus$ \ce{PbSO4(s)} (s) / \ce{Pb (s)}, $\oplus$ \ce{PbO2} (s) / \ce{PbSO4 (s)}, électrolyte aqueux \ce{H2SO4}

\solution{ 
\begin{minipage}{0.48\textwidth}
\begin{align*}
\ce{Pb (s) + H2SO4 (aq) &= PbSO4 (s) + 2 e^- + 2H^+ (aq)} \\
\ce{PbO2(s) + H2SO4(aq) + 2e^- + 2H^+ (aq) &= PbSO4(s) + 2 H2O(l)} \\
\hline
\ce{PbO2 (s) + Pb (s) + 2 H2SO4 (aq) &= 2PbSO4 (s) + 2 H2O (l)}
\end{align*}
\end{minipage}\begin{minipage}{0.48\textwidth}
\begin{center}
- %Réaction de formation (de \ce{PbSO4})
\end{center}
\end{minipage}}

\Q Pile Leclanché, pile saline : $\ominus$ \ce{Zn^{2+}(aq)/Zn(s)}, $\oplus$ \ce{MnO2(s)/MnOOH(s)}, électrolyte aqueux \ce{NH4Cl}

\solution{ 
\begin{minipage}{0.48\textwidth}
\begin{align*}
\ce{Zn(s) &= Zn^{2+} (aq) + 2e^-} \\
\ce{MnO2(s) + e^- + H^+ (aq) &= MnO(OH)(s)} \\
\hline
\ce{Zn(s) + 2 MnO2(s) + 2H^+ (aq) &= Zn^{2+} (aq) + 2 MnO(OH)(s)}
\end{align*}
\end{minipage}\begin{minipage}{0.48\textwidth}
\begin{center}
-
\end{center}
\end{minipage}}

\Q Pile alcaline : $\ominus$ \ce{Zn(OH)2/Zn}, $\oplus$ \ce{MnO2 / MnOOH}, électrolyte aqueux basique

\solution{ 
\begin{minipage}{0.48\textwidth}
\begin{align*}
\ce{Zn(s) + 2HO^-(aq) &= Zn(OH)2(s) + 2e^-} \\
\ce{MnO2(s) + e^- + H2O(l) &= MnO(OH)(s) + HO^-(aq)} \\
\hline
\ce{Zn(s) + 2MnO2(s) + 2H2O(l) &= Zn(OH)2(s) + 2MnO(OH)(s)}
\end{align*}
\end{minipage}\begin{minipage}{0.48\textwidth}
\begin{center}
-
\end{center}
\end{minipage}}

\Q Nickel-Cadmium : $\ominus$ \ce{Cd(OH)2 / Cd}, $\oplus$ \ce{NiO(OH)/Ni(OH)2}, électrolyte aqueux basique

\solution{ 
\begin{minipage}{0.48\textwidth}
\begin{align*}
\ce{Cd(s) + 2 HO^-(aq) &= Cd(OH)2(s) + 2e^-} \\
\ce{NiO(OH)(s) + e^- + H2O(l) &= Ni(OH)2(s) + HO^-(aq)} \\
\hline
\ce{Cd (s) + 2NiO(OH)(s) + 2H2O(l) &= Cd(OH)2(s) + Ni(OH)2(s)}
\end{align*}
\end{minipage}\begin{minipage}{0.48\textwidth}
\begin{center}
-
\end{center}
\end{minipage}}

\Q Nickel-métal-hydrure : $\ominus$ \ce{M/MH} , $\oplus$ \ce{NiO(OH)/Ni(OH)2}, électrolyte aqueux basique

\solution{ 
\begin{minipage}{0.48\textwidth}
\begin{align*}
\ce{MH (s) + HO^- (aq) &= M(s) + H2O(l) + e^-} \\
\ce{NiO(OH)(s) + e^- + H2O(l) &= Ni(OH)2(s) + HO^-(aq)} \\
\hline
\ce{NiO(OH) (s) + MH(s) &= M(s) + 2 Ni(OH)2(s)}
\end{align*}
\end{minipage}\begin{minipage}{0.48\textwidth}
\begin{center}
Réaction d'insertion (de l'hydrure \ce{H^-})
\end{center}
\end{minipage}}